\scanpage{7}
\uline{Prop. 2}

Soit $G$ un schéma en groupes sur $S$ loc noeth,
on suppose que $\text{\uuline{Hilb}}_{\inserted{G/S}}\deleted{(G/S)}$ existe. On s'intéresse
% slanted
\margin{$G$ \unsure{qu.|q.} projectif sur $S$ ???}\unsure{à ens|à l'ens|l'ens|à l'ens.} $F(S)\subset\mathrm{Hilb}(G/S)$ formé des
sous-schémas fermés de $G$, propres et plats sur $S$,
qui sont des \uline{schémas \unsure{de|en} ss-groupes}. Pour tout
$S'$, sur $S$, on définit de $=$ $F(S')\subset\mathrm{Hilb}(G'/S')$.
d'où un ``sous foncteur'' de $\text{\uuline{Hilb}}_{G/S}$.

\uline{Je dis qu'il est représentable} et que
\uline{le préschéma} $M$ \uline{sur} $S$ \uline{est tel que le
morphisme induit} $M\mathrel{-}\text{\uuline{Hilb}}_{G/S}$ \uline{est
une immersion fermée}. \uline{Si} $G$ est
\uline{propre sur} $S$ \uline{et} $\mathrm{Hilb}_{G/S}$ \uline{est réunion} disjointe
\uline{d'ouverts} \deleted{\unsure{qui} \unsure{pr}}\inserted{$H_i$} \uline{de type fini \deleted{\unsure{$H_i$}} sur} $S$,
\deleted{\unsure{alors} les $H_i\cap M=M_i$} \inserted{donc} $M$ \unsure{--|est}
\uline{réunion disjointe des ouverts} $M_i=M\cap H_i$,
\uline{alors les} $M_i$ \uline{sont propres sur} $S$,

N.B. Si on s'intéresse à la \unsure{classification|caractérisation} des \inserted{ss} schémas \deleted{\unsure{s}}\inserted{en \unsure{grpes|groupes}} $H$
de $G$ qui sont \inserted{de \uline{plus} (déjà propres !)} finis (\unsure{resp$^t$|resp.} séparables, resp$^t$ abéliens),
on trouve des ss-schémas \uline{ouverts} de $M$, les premiers
\unsure{étant} également \uline{fermés}.
% archive sheet number: 6

